\section{Penelitian Terkait}
Penelitian yang dilakukan oleh Al-Ismail et al. \cite{7} membahas penerapan teknologi \textit{blockchain} dalam ekosistem Sistem Informasi Rumah Sakit (\textit{Hospital Information System} - HIS). Latar belakang utama penelitian ini adalah kerentanan pada sistem manajemen data medis terpusat tradisional terhadap manipulasi data, kebocoran privasi, dan akses tidak sah.

Dalam penelitian tersebut, diusulkan sebuah kerangka kerja manajemen data rekam medis berbasis \textit{blockchain} yang bertujuan untuk meningkatkan integritas, ketersediaan (\textit{availability}), dan nir-penyangkalan (\textit{non-repudiation}) data. Solusi ini memanfaatkan sifat desentralisasi dan \textit{immutability} dari \textit{blockchain} untuk memberikan kendali yang lebih besar kepada pasien atas data kesehatan mereka, serta memfasilitasi interoperabilitas yang aman antar penyedia layanan kesehatan. Hasil penelitian menunjukkan bahwa penggunaan \textit{smart contract} dan buku besar terdistribusi mampu meminimalkan risiko penipuan penagihan (\textit{fraudulent billing}) serta menjamin bahwa riwayat medis pasien tersimpan secara akurat dan tidak dapat diubah oleh pihak ketiga tanpa otorisasi.

Meskipun penelitian Al-Ismail et al. berhasil membuktikan efektivitas \textit{blockchain} dalam aspek keamanan preventif dan integritas data, penelitian tersebut belum secara spesifik membahas mekanisme penanganan insiden (\textit{incident response}) ketika serangan tetap terjadi, khususnya yang berasal dari aktor internal (\textit{insider threats}) dengan hak akses sah. Selain itu, aspek manajemen \textit{audit trail} yang komprehensif untuk kebutuhan forensik digital juga belum menjadi fokus utama. Oleh karena itu, penelitian Tugas Akhir ini akan melengkapi kekurangan tersebut dengan mengintegrasikan sistem deteksi dan respons insiden menggunakan ELK Stack sebagai SIEM serta penyimpanan log forensik \textit{off-chain} pada IPFS, untuk menjamin akuntabilitas sistem secara menyeluruh.